\documentclass[dvipdfmx,b5j,10pt]{jsarticle}%

\usepackage{mypackage,chuushaku}%

\begin{document}%

\chuukigou{}{\rotatebox[origin=c]{180}{\textcolor{gray}{\▲}}}
\begin{ascolorbox4}{問題}\parindent=2zw
I have a good memory, except for names and faces, and I do not forget what I have read. The disadvantage of this is that having read all the great novels of the world two or three times I can no longer read them with relish. There are few modern novels that excite my interest, and I do not know what I should do for relaxation were it not for the innumerable detective stories that so engagingly pass the time and once read pass straight out of one's mind.\source{Somerset Maugham, \it{A Writer's Notebook}}
\end{ascolorbox4}\br{-.5}%
\begin{chuushaku}[.7]<.5>
\begin{鉄壁カラーボックス}{\!\!表現型\!\!}\noindent\mc{随筆文逸話文体}\end{鉄壁カラーボックス}%
\br0~\chuu{私は記憶力が良い}\br{-1}\sankakuline*{I have a good memory}%
SVOを直訳してもよいが，\point{名詞構文の訳し方}\noindent に従って，I memorize (things) well「物を上手に記憶する；記憶力がよい」と言い換えてもよい。%
\br0~\chuu{人の名前や顔は別として}\br{-1}\sankakuline*{except for names and faces}%
「名前と顔以外は」でよいが，作者の意図としては「人の名前と顔を覚えるのは\kenten{苦手}だ」という含蓄がある。%
\br0~\chuu{一度読んだ本は忘れることがありません}\br{-1}\sankakuline*{I do not forget what I have read.}%
\point{wh-節内の現在完了の訳し方}\point{関係詞の訳し方}\noindent を踏まえ，「\kenten{一度}読み終えた\kenten{本}」とすると自然。%
\br0~\chuu{そのため不都合なことに／\\そのため不便な点は}\br{-1}\sankakuline*{The disadvantage of this is that}%
\point{同格名詞節の訳し方}%
\noindent に則ると，「不都合なことに」と訳せる。もちろん「不便な点は」でも問題ない。%
\end{chuushaku}

\end{document}％